% ═══════════════════════════════════════════════════════════════════════════
% ÜBERSICHT DS-03: Modalverben (poder, querer, tener que)
% Handout für SuS (A4 Hochformat)
% ═══════════════════════════════════════════════════════════════════════════

\documentclass[11pt, a4paper]{article}

\usepackage[utf8]{inputenc}
\usepackage[T1]{fontenc}
\usepackage[ngerman]{babel}
\usepackage{helvet}
\renewcommand{\familydefault}{\sfdefault}

\usepackage[margin=1.5cm, top=1.2cm, bottom=1.2cm]{geometry}
\usepackage[table]{xcolor}
\usepackage{array}
\usepackage{tabularx}
\usepackage{booktabs}
\usepackage{tikz}
\usepackage{tcolorbox}
\tcbuselibrary{skins}

% Farben
\definecolor{spanisch}{HTML}{c41e3a}
\definecolor{poder}{HTML}{1565c0}
\definecolor{querer}{HTML}{2a7a4b}
\definecolor{tenerque}{HTML}{b35c00}
\definecolor{hellgrau}{HTML}{f5f5f5}

\pagestyle{empty}

\begin{document}

% ═══════════════════════════════════════════════════════════════════════════
% TITEL
% ═══════════════════════════════════════════════════════════════════════════

\begin{center}
\begin{tikzpicture}
\fill[spanisch] (0,0) rectangle (18,0.9);
\node[white, font=\Large\bfseries] at (9,0.45) {Los verbos modales -- Modalverben};
\end{tikzpicture}
\end{center}

\vspace{0.3cm}

\begin{center}
\textit{Modalverben + Infinitiv = Was man kann, will oder muss}
\end{center}

\vspace{0.3cm}

% ═══════════════════════════════════════════════════════════════════════════
% PODER
% ═══════════════════════════════════════════════════════════════════════════

\begin{tcolorbox}[
    colback=poder!10,
    colframe=poder,
    title={\color{white}\bfseries PODER -- können},
    fonttitle=\bfseries
]

\begin{tabular}{@{}p{6cm}p{10.5cm}@{}}

\textbf{Konjugation:}\\[0.2cm]
\begin{tabular}{ll}
yo & \textbf{puedo} \\
tú & \textbf{puedes} \\
él/ella & \textbf{puede} \\
nosotros & \textbf{podemos} \\
vosotros & \textbf{podéis} \\
ellos & \textbf{pueden} \\
\end{tabular}

\footnotesize\textit{(Stammvokalwechsel: o $\rightarrow$ ue)}

&

\textbf{Verwendung:} Fähigkeit, Möglichkeit, Erlaubnis

\vspace{0.2cm}

\textbf{Beispiele:}
\begin{itemize}
    \item \textit{Puedo hablar español.} -- Ich kann Spanisch sprechen.
    \item \textit{¿Puedes ayudarme?} -- Kannst du mir helfen?
    \item \textit{No podemos ir al cine.} -- Wir können nicht ins Kino gehen.
\end{itemize}

\vspace{0.1cm}
\footnotesize\textbf{Struktur:} poder (konjugiert) + Infinitiv

\\
\end{tabular}

\end{tcolorbox}

\vspace{0.3cm}

% ═══════════════════════════════════════════════════════════════════════════
% QUERER
% ═══════════════════════════════════════════════════════════════════════════

\begin{tcolorbox}[
    colback=querer!10,
    colframe=querer,
    title={\color{white}\bfseries QUERER -- wollen / mögen},
    fonttitle=\bfseries
]

\begin{tabular}{@{}p{6cm}p{10.5cm}@{}}

\textbf{Konjugation:}\\[0.2cm]
\begin{tabular}{ll}
yo & \textbf{quiero} \\
tú & \textbf{quieres} \\
él/ella & \textbf{quiere} \\
nosotros & \textbf{queremos} \\
vosotros & \textbf{queréis} \\
ellos & \textbf{quieren} \\
\end{tabular}

\footnotesize\textit{(Stammvokalwechsel: e $\rightarrow$ ie)}

&

\textbf{Verwendung:} Wunsch, Wille, Vorliebe

\vspace{0.2cm}

\textbf{Beispiele:}
\begin{itemize}
    \item \textit{Quiero ser médico.} -- Ich will Arzt werden.
    \item \textit{¿Quieres un café?} -- Möchtest du einen Kaffee?
    \item \textit{Queremos viajar a España.} -- Wir wollen nach Spanien reisen.
\end{itemize}

\vspace{0.1cm}
\footnotesize\textbf{Struktur:} querer (konjugiert) + Infinitiv / Nomen

\\
\end{tabular}

\end{tcolorbox}

\vspace{0.3cm}

% ═══════════════════════════════════════════════════════════════════════════
% TENER QUE
% ═══════════════════════════════════════════════════════════════════════════

\begin{tcolorbox}[
    colback=tenerque!10,
    colframe=tenerque,
    title={\color{white}\bfseries TENER QUE -- müssen},
    fonttitle=\bfseries
]

\begin{tabular}{@{}p{6cm}p{10.5cm}@{}}

\textbf{Konjugation:}\\[0.2cm]
\begin{tabular}{ll}
yo & \textbf{tengo que} \\
tú & \textbf{tienes que} \\
él/ella & \textbf{tiene que} \\
nosotros & \textbf{tenemos que} \\
vosotros & \textbf{tenéis que} \\
ellos & \textbf{tienen que} \\
\end{tabular}

\footnotesize\textit{(1. Pers. unregelmäßig + Stammvokal e $\rightarrow$ ie)}

&

\textbf{Verwendung:} Notwendigkeit, Pflicht, Zwang

\vspace{0.2cm}

\textbf{Beispiele:}
\begin{itemize}
    \item \textit{Tengo que estudiar.} -- Ich muss lernen.
    \item \textit{Tienes que hacer los deberes.} -- Du musst die Hausaufgaben machen.
    \item \textit{Tenemos que irnos.} -- Wir müssen gehen.
\end{itemize}

\vspace{0.1cm}
\footnotesize\textbf{Struktur:} tener (konjugiert) + \textbf{que} + Infinitiv

\\
\end{tabular}

\end{tcolorbox}

\vspace{0.4cm}

% ═══════════════════════════════════════════════════════════════════════════
% ZUSAMMENFASSUNG
% ═══════════════════════════════════════════════════════════════════════════

\begin{center}
\fcolorbox{spanisch}{hellgrau}{%
\begin{minipage}{0.95\textwidth}
\centering
\textbf{Zusammenfassung: Modalverb (konjugiert) + Infinitiv}

\vspace{0.3cm}

\begin{tabular}{|c|c|c|c|}
\hline
\rowcolor{spanisch!20}
\textbf{Verb} & \textbf{Bedeutung} & \textbf{Stammwechsel} & \textbf{Beispiel} \\
\hline
\textcolor{poder}{\textbf{poder}} & können & o $\rightarrow$ ue & \textit{Puedo nadar.} \\
\hline
\textcolor{querer}{\textbf{querer}} & wollen & e $\rightarrow$ ie & \textit{Quiero comer.} \\
\hline
\textcolor{tenerque}{\textbf{tener que}} & müssen & e $\rightarrow$ ie (+ \textbf{que}!) & \textit{Tengo que ir.} \\
\hline
\end{tabular}

\end{minipage}%
}
\end{center}

\end{document}
